\section{Термодинамика микроуровня}

\begin{axiom}[A5 · Третье начало (Нернст--Планк)]
\label{ax:A5}
Третье начало термодинамики: при $T\to 0$ энтропия равновесной системы стремится к конечному
пределу (Нернст); в соглашении Планка
\[
  S \xrightarrow[T\to 0]{} 0.
\]
На~$\layerMicro$ при $z=0$ фазовый поворот $z e^{i\Phi}$ обнуляет поле
и останавливает эволюцию.
При $|z|\to 0$ фазовый вращатель~$\Phi$ неограниченно возрастает: вакуум --- \emph{осциллирующий} фон
с $\rho_{\mathrm{field}}>0$.
\end{axiom}

\begin{axiom}[A6 · Второй закон (локально)]
\label{ax:A6}
На каждой $\varepsilon$-окрестности локальная фазовая энтропия (при фиксированной норме $N$)
не убывает: разрешено перемешивание / перераспределение фаз, запрещена «заморозка'' всех фаз в rest.
Рост энтропии не отменяет A3 --- не означает $N\to 0$.
\end{axiom}

\begin{axiom}[A7 · Планковский потолок плотности]
\label{ax:A7}
Энергетическая плотность ограничена планковским потолком $u_P$:
\[
  \rho_E(x,t)=\mu_P c^2 |z(x,t)|^2 \le u_P
  \quad\Leftrightarrow\quad
  |z|^2 \le 1 \ \text{(natural)}.
\]
Насыщающий вращатель использует знаменатель $+\varepsilon$, $\varepsilon\leftrightarrow u_P$.
\end{axiom}

\begin{axiom}[A8 · Макро-линейность]
\label{ax:A8}
При больших $|z|$ нелинейность вращателя затухает --- на $\Tlayer$ поле становится близким к линейному
(гладкость макроосреднение).
\end{axiom}

\begin{figure}[htbp]
  \centering
  \begin{tikzpicture}[carrier panel]
  \draw[carrier object, domain=0:1.15, samples=60] plot (\x, {(\x-0.35)^2+0.05});
  \draw[carrier muted, dashed] (0.35,0) -- (0.35,1.0);
  \node at (0,1.05) {$\times$};
  \node[anchor=west, font=\scriptsize] at (0.05,0.85) {deadlock};
  \fill[carrier object] (0.35,0.05) circle (1.5pt);
  \node[anchor=west, font=\scriptsize] at (0.48,0.18) {вакуум $z_{\min}>0$};
  \node[carrier muted, font=\scriptsize] at (0.55,-0.45) {$|z|$};
  \node[carrier muted, font=\scriptsize, rotate=90] at (-0.25,0.55) {отклик $\symCapitalPhi$};
\end{tikzpicture}

  \caption{$A_5$: недостижимый ноль амплитуды, осциллирующий вакуум $z_{\min}>0$.}
  \label{fig:axiom-a5-vacuum}
\end{figure}

\begin{figure}[htbp]
  \centering
  \begin{tikzpicture}[carrier panel]
  \foreach \i in {0,...,11} {
    \pgfmathsetmacro{\ph}{30*\i}
    \pgfmathsetmacro{\x}{1.05*cos(\ph)}
    \pgfmathsetmacro{\y}{1.05*sin(\ph)}
    \draw[carrier object, carrier fill]
      (\x-0.08,\y-0.22) rectangle (\x+0.08,\y+0.22);
  }
  \draw[carrier object, dashed] (-1.1,0) -- (1.1,0);
  \node[carrier muted, font=\scriptsize] at (0.55,0.75) {$N=\sum|z|^2$};
\end{tikzpicture}

  \caption{$A_6$: рост фазовой энтропии при фиксированном $N$.}
  \label{fig:axiom-a6-entropy}
\end{figure}

